\documentclass[UTF8, a4paper, 12pt]{ctexart}

\usepackage{amsmath, amssymb, amsthm}
\usepackage{bm}
\usepackage{geometry}
\usepackage{graphicx}
\usepackage{booktabs}
\usepackage{enumitem}
\usepackage{hyperref}
\usepackage{listings}
\usepackage{xcolor}

\geometry{margin=2.5cm}
\hypersetup{
  colorlinks=true,
  linkcolor=black,
  urlcolor=blue,
  citecolor=black
}

\lstset{
  language=C++,
  basicstyle=\small\ttfamily,
  keywordstyle=\color{blue!60!black}\bfseries,
  commentstyle=\color{green!40!black},
  stringstyle=\color{purple},
  breaklines=true,
  frame=single,
  columns=fullflexible,
  tabsize=2
}

\theoremstyle{definition}
\newtheorem{problem}{题目}
\theoremstyle{plain}
\newtheorem{proposition}{命题}
\newtheorem{lemma}{引理}

\title{偏微分方程作业 1}
\author{姓名：\underline{\hspace{4cm}} \quad 学号：\underline{\hspace{4cm}}}
\date{\today}

\begin{document}
\maketitle

\tableofcontents
\vspace{1em}

\section{作业说明}

本次作业的材料分三处存放：

\begin{itemize}[leftmargin=2em]
  \item 本文件：解析推导与文字解答；
  \item \texttt{../cpp/}：C++ 数值实现；
  \item \texttt{../viz/}：HTML 可视化。
\end{itemize}

编译本笔记：

\begin{verbatim}
latexmk -xelatex notes/main.tex
\end{verbatim}

\section{题目与解答}

\begin{problem}[请替换为实际题目]
在区间 \(x\in[0,1]\)、时间 \(t\ge 0\) 上考虑一维热方程
\begin{equation}
  u_t = k u_{xx},
  \qquad
  u(0,t)=u(1,t)=0,
  \qquad
  u(x,0)=f(x).
\end{equation}
讨论适定性，并给出一种稳定的有限差分格式。
\end{problem}

\subsection*{解答}

将题目、推导、中间步骤写在这里。需要插图时，把文件放到 \texttt{figures/}：

\begin{verbatim}
\includegraphics[width=0.8\linewidth]{figures/example.png}
\end{verbatim}

\section{数值方法}

仓库中的示例程序求解一维热方程的显式差分（FTCS）：
\begin{equation}
  \frac{u_j^{n+1}-u_j^n}{\Delta t}
  =
  k\frac{u_{j+1}^n-2u_j^n+u_{j-1}^n}{(\Delta x)^2}.
\end{equation}
稳定条件为
\begin{equation}
  r = \frac{k\Delta t}{(\Delta x)^2} \le \tfrac12.
\end{equation}

默认初值取中间的高斯包络
\[
  f(x)=\exp\bigl(-80(x-1/2)^2\bigr),
\]
两端取齐次 Dirichlet 条件。若 \(r>1/2\)，程序会自动加大 \texttt{nt} 以满足稳定条件。按实际题目修改 \texttt{cpp/src/heat.cpp} 即可。

\section{程序与可视化}

\begin{verbatim}
make cpp    # 写出 data/solution.json
make viz    # 浏览器打开 http://localhost:8000/viz/
\end{verbatim}

运行后把关键截图贴到本节，并简要说明网格、时间步与观察到的现象。

\section{参考}

\begin{enumerate}[leftmargin=2em]
  \item 课程讲义与课堂笔记。
  \item Evans, L. C. \textit{Partial Differential Equations}.
\end{enumerate}

\end{document}
